\section{Discussion}
\label{sec:conclusion}

\textbf{Conclusion.} We have developed \benchmark{}, a novel lightweight framework that facilitates interaction between Language Models and the underlying environment, enabling them to mimic the human approach to language-to-code generation. Our framework has shown promising results when applied to state-of-the-art models using different prompting styles. It effectively leverages the capabilities of LMs to break down complex tasks and recover from errors within a secure and isolated environment. The ability to seamlessly convert existing datasets into the interactive format using \texttt{\benchmark{}Env} API, and furthermore, the Bash and SQL environments, empowers task designers to construct new tasks to unlock the plethora of challenges that await in the space of interactive coding.

\textbf{Limitations and future directions.} We point out several current limitations of \benchmark{}. At this time, the number of \benchmark{} based environments is limited to Bash, SQL, and Python action spaces and datasets; within the near future, we plan to expand the number of offerings to cover a wider set of programming languages and datasets that should further deliver on \benchmark{}'s purported promises of efficient and expressive task construction. Second, the CTF dataset is limited to just four task instances due to our manual curation procedure. We hope to release more formal work soon that provides a more thorough analysis of the reasoning and collaboration challenges of the CTF task along with a more extensive dataset for evaluation purposes.